\documentclass[11pt]{article}

\usepackage[T1]{fontenc}
\usepackage[utf8]{inputenc}
\usepackage[english]{babel}
\usepackage[margin=1in]{geometry}
\usepackage{lmodern}
\usepackage{microtype}
\usepackage{mathtools}
\usepackage{amsmath,amssymb,amsthm}
\usepackage{array}
\usepackage{booktabs}
\usepackage{tabularx}
\usepackage{float}
\usepackage{graphicx}
\usepackage{caption}
\usepackage{enumitem}
\usepackage{fancyhdr}
\usepackage[numbers]{natbib}
\usepackage{doi}
\usepackage{xurl}
\usepackage{hyperref}
\usepackage{cleveref}
\usepackage{orcidlink}
\usepackage{titling}

\newcommand{\BranchwiseQ}{K}
\newcommand{\RecombinationQ}{R}
\newcommand{\ForgetfulMap}{\eta}
\newcommand{\BranchwiseEq}{\sim_{\mathrm{bw}}}
\newcommand{\RecombinationEq}{\sim_{\mathrm{rec}}}
\newcommand{\ObservableFamily}{\mathcal{O}}
\newcommand{\Assemblage}{\mathcal{A}}
 
\newtheorem{theorem}{Theorem}[section]
\newtheorem{lemma}[theorem]{Lemma}
\newtheorem{proposition}[theorem]{Proposition}
\newtheorem{corollary}[theorem]{Corollary}
\theoremstyle{definition}
\newtheorem{definition}[theorem]{Definition}
\newtheorem{remark}[theorem]{Remark}
\floatstyle{ruled}
\newfloat{algorithm}{tbp}{loa}
\floatname{algorithm}{Algorithm}

\urlstyle{same}
\captionsetup{
  font=small,
  labelfont=bf,
  labelsep=period
}
\captionsetup[table]{skip=6pt,position=top}
\captionsetup[figure]{skip=6pt}
\captionsetup[algorithm]{skip=6pt,position=top}
\crefname{algorithm}{algorithm}{algorithms}
\Crefname{algorithm}{Algorithm}{Algorithms}
\crefname{theorem}{theorem}{theorems}
\Crefname{theorem}{Theorem}{Theorems}
\crefname{lemma}{lemma}{lemmas}
\Crefname{lemma}{Lemma}{Lemmas}
\crefname{proposition}{proposition}{propositions}
\Crefname{proposition}{Proposition}{Propositions}
\crefname{corollary}{corollary}{corollaries}
\Crefname{corollary}{Corollary}{Corollaries}
\crefname{definition}{definition}{definitions}
\Crefname{definition}{Definition}{Definitions}
\crefname{remark}{remark}{remarks}
\Crefname{remark}{Remark}{Remarks}
\setlist[enumerate]{leftmargin=*, itemsep=2pt, topsep=4pt}
\setlength{\droptitle}{-3.2em}
\pretitle{\begin{center}\LARGE\bfseries}
\posttitle{\par\end{center}\vspace{0.5em}}
\preauthor{\begin{center}\normalsize}
\postauthor{\par\end{center}\vspace{-0.4em}}
\predate{\begin{center}\small}
\postdate{\par\end{center}\vspace{0.6em}}
\setlength{\emergencystretch}{2em}
\clubpenalty=10000
\widowpenalty=10000
\displaywidowpenalty=10000
\interfootnotelinepenalty=10000
\allowdisplaybreaks[2]
\fancypagestyle{firstpage}{\fancyhf{}
  \fancyfoot[C]{\scriptsize Automorph Inc., Wilmington, DE, USA \quad \texttt{ioannis@automorph.io} \quad \textcopyright\ 2026 \quad CC-BY 4.0}
  \renewcommand{\headrulewidth}{0pt}
  \renewcommand{\footrulewidth}{0pt}
}

\title{Internal Structure of a Surviving Promoted Package on Fixed Support in Six Birds Theory}
\author{Ioannis Tsiokos\,\orcidlink{0009-0009-7659-5964}}
\date{April 2026}
\hypersetup{
  pdftitle={Internal Structure of a Surviving Promoted Package on Fixed Support in Six Birds Theory},
  pdfauthor={Ioannis Tsiokos},
  hidelinks
}

\begin{document}

\maketitle
\thispagestyle{firstpage}

\begin{abstract}
Six Birds Theory constructs layered descriptions of a stochastic substrate without importing Hilbert space or the Born rule. Earlier papers in the series showed that the substrate supports interference-like recombination structure on fixed support and that at least one promoted package survives on the Route~A flagship family, but left open the internal status of that survivor. We classify it here on a frozen exact-finite corpus consisting of two flagship anchors, four nearby positives, matched controls, and one support-only comparison row, comparing intrinsic, sectioned, bundle-section-fiber, and protocol-indexed candidates under descended readout, signature-level descended update, second-stage recoarsening, protocol honesty, invariant validation, and exact-finite quadratic probes. The verdict is an intrinsic package: on the flagship core and required controls the intrinsic nonindexed candidate is the only candidate that combines descended readout, a nontrivial second-stage stability plateau, retained invariant and quadratic families, and protocol honesty. The nearby positives remain anchor-only on the frozen exact readout surface, but carrier diagnostics and family-level quotient evidence support an erased-partner interpretation. Weak primitive-action participation and a flagship-local exact chart are recorded as downstream, non-gating diagnostics.
\end{abstract}

\medskip
\noindent\textbf{Keywords:} fixed support; promoted package; descended structure; recoarsening stability; exact-finite quadratic probes; quotient structure.

\section{Introduction}

Six Birds Theory is an autonomous emergence framework that constructs layered descriptions of a stochastic substrate from stable records, package-relative identity, and admissible completion, without importing Hilbert space, amplitudes, or the Born rule \cite{Tsiokos2026EventPackages,Tsiokos2026PredictiveQuotients}. A central question for the theory is whether its substrate supports not only interference-like phenomena at the recombination level but also genuine higher-level objects---promoted packages with their own internal structure---that are not merely compressed summaries of the lower layer.

Earlier papers in the series developed the fixed-support setting in stages. The first established admissible event packages and the discipline of comparing candidate descriptions on a common support \cite{Tsiokos2026EventPackages}. The second introduced current and predictive quotients on single histories and showed that hidden predictive structure can exist below the current-event interface \cite{Tsiokos2026PredictiveQuotients}. The third lifted the same discipline to branch assemblages and showed that branchwise predictive content and later recombination-sensitive content can separate on fixed support---the substrate-level analogue of which-path interference \cite{Tsiokos2026MarkingErasure}. The fourth isolated a control-stable Route~A flagship family with a theorem-backed Lean kernel \cite{Tsiokos2026RecombinationWitnesses}. The fifth asked whether that witness family supports a genuine promoted package and showed that at least one intrinsic candidate survives on fixed support, but did not determine what kind of promoted object had survived \cite{Tsiokos2026PromotionCriteria}.

The question addressed here is therefore internal rather than existential. Is the survivor an intrinsic promoted package with its own descended structure? A section-dependent chart on a deeper object? Protocol-indexed, or merely a useful lower-layer compression? The answer matters for the broader program: if the substrate supports a promoted object with genuine internal organization, then the layered structure built by the series is not merely a bookkeeping convenience but a real feature of the emergence framework. Answering requires internal structure that survives the relevant honesty and stability checks on fixed support, not an assumed carrier.

We study a frozen exact-finite corpus consisting of two flagship anchors, four nearby positives, four required controls, and one support-only comparison row. On that corpus we compare four candidate object types inherited from the previous paper: intrinsic nonindexed, sectioned, bundle-section-fiber, and protocol-indexed. The audit proceeds through descended readout, signature-level descended update, admissible second-stage recoarsening, protocol honesty, invariant validation, exact-finite quadratic probes, and a separate nearby-positive synthesis layer.

\paragraph{Contributions.}
First, we show that the best-supported verdict on the frozen corpus is an intrinsic package: the intrinsic nonindexed candidate is the only candidate that combines descended readout, signature-level descended update support, and a nontrivial second-stage stability plateau on the flagship core and required controls. Second, we validate retained invariant families and exact-finite quadratic families on that intrinsic plateau, separating claim-bearing families from comparison-only and proxy-only quantities. Third, we show that sectioned and bundle-section-fiber alternatives remain informative comparison branches but do not displace the intrinsic verdict, and that dual augmentation adds no decisive improvement. Fourth, we close the nearby-positive route with a caveat: on the frozen exact readout surface the nearby positives remain anchor-only, yet detector-independent carrier diagnostics and family-level quotient evidence support an erased-partner interpretation under a restricted detector-unlock regime.

\paragraph{Relation to existing mathematics.}
The mathematical tools used in this paper are individually standard; the contribution is their composition into an internal-structure audit on the Six Birds substrate and the demonstration that the audit yields positive results. Descended readout and update are instances of the quotient-descent and aggregation questions studied in behavioral equivalence and exact lumpability \cite{Sangiorgi2009,VanDerSchaft2004,KemenySnell1976,Buchholz1994}. Second-stage stability under recoarsening is a form of the minimality and refinement questions that appear in model minimization for stochastic systems \cite{GivanDeanGreig2003}. Protocol honesty---requiring that apparent structure not depend on omitted protocol state---extends the internalization discipline already used in the series. The Lean formal layer uses standard machine-checked finite mathematics \cite{deMouraUllrich2021,MathlibCommunity2020}. As in the earlier papers, the contribution is not the introduction of these general mathematical ideas but their application, under exact-finite control and matched comparisons, to a specific emergence substrate \cite{ShaliziCrutchfield2001}.

\paragraph{Scope.}
We classify the internal status of the surviving promoted object on the Route~A flagship family and its fixed comparison surface. We do not derive any ambient carrier model, do not claim a readout-surface extension from the anchors to the nearby positives, do not promote bundle, dual, or indexed alternatives to the main verdict, and do not certify a fully autonomous higher-level layer. The flagship-local chart reported later is local and descriptive only. The nearby-positive synthesis does not rewrite the earlier readout-surface limitation; it supplements it with a separate family-level interpretation under a restricted detector-unlock caveat. Downstream primitive-participation and local-chart results are reported but remain non-gating.

\paragraph{Roadmap.}
Section 2 recalls the fixed-support Route A starting surface and the precise boundary left by the previous paper. Section 3 states the internal-structure criteria and theorem spine. Section 4 fixes the canonical corpus, candidate classes, and audit pipeline. Sections 5 and 6 establish descended structure and second-stage stability, leading to the intrinsic verdict. Section 7 develops the retained invariant and exact-finite quadratic families. Section 8 explains the nearby-positive deadlock and its detector-limited erased-partner resolution. Section 9 records downstream diagnostics on weak primitive participation and flagship-local charts. Section 10 summarizes the formal criterion layer. Section 11 discusses limitations, and Section 12 concludes.
 \section{Fixed-Support Background and Route A Starting Surface}

\subsection{Fixed support across the series}

Six Birds Theory constructs its descriptions from a finite stochastic substrate: stable records, package-relative identity, and admissible completion provide the working materials, while fixed support provides the comparison discipline. That discipline is designed to prevent apparent structural gains from being purchased by silently changing the underlying comparison base: the relevant question is always what survives, obstructs, or refines on a common fixed support. Within that regime, the series has developed the event-package setting, the predictive quotient layer beneath the current interface, recombination-sensitive witness structure, and the control-stable Route~A flagship family on which promoted-package questions can be asked sharply \cite{Tsiokos2026EventPackages,Tsiokos2026PredictiveQuotients,Tsiokos2026MarkingErasure,Tsiokos2026RecombinationWitnesses,Tsiokos2026PromotionCriteria}.

At a broader conceptual level, the local-to-global obstruction language used here is close in spirit to sheaf-theoretic contextuality, while the predictive quotient layer belongs to the wider family of predictive-state and causal-state abstractions \cite{AbramskyBrandenburger2011,ShaliziCrutchfield2001,LittmanSuttonSingh2001}.

The progression of those papers matters for the present starting point. The first paper established the fixed-support event-package framework and the admissible-global-packaging viewpoint. The second introduced predictive quotient structure below the current interface. The third isolated recombination-sensitive exact-finite witnesses. The fourth narrowed that machinery to the Route~A flagship family and its control-stable anchors. The fifth then asked the existential promoted-package question on that frozen surface and answered it positively, but deliberately stopped short of identifying the internal status of the survivor.

\subsection{What the previous paper settled}

More precisely, the previous paper settled the survival question on a frozen Route~A corpus. On that surface, a promoted package survives on fixed support; the intrinsic route is the primary survivor; lift- or representative-selected routes remain real but secondary; and slice/shadow or defect-style surfaces remain support-only or explanatory. What was not settled was the internal classification: whether the survivor is best understood as an intrinsic promoted package with its own descended structure, a section-dependent chart, a protocol-indexed object, a bundle-like enhancement, or a lower-layer compression \cite{Tsiokos2026PromotionCriteria}.

\subsection{Inherited frozen surface}

The present paper therefore starts from a frozen inherited surface rather than a new search. At the manuscript level, that surface is decomposed as
\[
\mathcal C
  = \mathcal C_{\mathrm{flag}}
    \sqcup \mathcal C_{\mathrm{near}}
    \sqcup \mathcal C_{\mathrm{ctrl}}
    \sqcup \mathcal C_{\mathrm{supp}},
\qquad
\lvert \mathcal C_{\mathrm{flag}} \rvert = 2,\ 
\lvert \mathcal C_{\mathrm{near}} \rvert = 4,\ 
\lvert \mathcal C_{\mathrm{ctrl}} \rvert = 4,\ 
\lvert \mathcal C_{\mathrm{supp}} \rvert = 1.
\]
The corpus is fixed in advance and is not enlarged here. The present work is not a route-zoo expansion or a new discovery pass, but an internal-structure audit on the exact surface that the previous paper left behind.

\subsection{Open question}

The key open issue is internal structure. A surviving promoted package is more than a convenient label only if it supports its own descended readout and update structure, remains honest under protocol treatment, retains nontrivial content under admissible second-stage recoarsening, and carries invariant or quadratic content that is not merely proxy bookkeeping from below. Those are the issues taken up in the sections that follow.
 \section{Internal-Structure Criteria and Theorem Spine}

\subsection{Candidate packages and descended interfaces}

The present paper compares a small candidate family rather than assuming a carrier in advance. The candidates are fourfold: intrinsic, sectioned, bundle-section-fiber, and protocol-indexed \cite{Tsiokos2026PromotionCriteria}. They are hypotheses about internal object type, not implementation conveniences.

\begin{definition}[Candidate promoted package]
\label{def:candidate-package}
A \emph{candidate promoted package} on a fixed corpus consists of a candidate kind together with the promoted-state data and admissible representational surface carried by that kind. The intrinsic candidate carries only the promoted base state, the sectioned candidate carries the base state together with an admissible section choice, the bundle-section-fiber candidate carries the base state together with an admissible section fiber, and the protocol-indexed candidate carries the base state together with explicit protocol or index data.
\end{definition}

\begin{definition}[Descended readout and signature-level descended update]
\label{def:descended-interfaces}
Let $P$ be a candidate promoted package on a fixed comparison surface $\mathcal C$. We say that $P$ has \emph{descended readout support} on $\mathcal C$ if the designated promoted-level readout family descends honestly from the lower layer on that surface. We say that $P$ has \emph{signature-level descended update support} on $\mathcal C$ if the retained lower-generator interface descends to a coherent promoted-level action or closure interface on $\mathcal C$. In this paper, the update interface is signature-level rather than a full executable promoted dynamics.
\end{definition}

\subsection{Joint criterion for internal structure}

Internal structure is not certified by one scalar statistic. A promoted object counts as internally real only if several conditions are met jointly: descended readout, signature-level descended update, admissible second-stage stability on its retained surface, protocol honesty, validated invariant content, and retained exact-finite quadratic content.

\begin{definition}[Second-stage stability and protocol honesty]
\label{def:stability-honesty}
Let $P$ be a candidate promoted package. We say that $P$ is \emph{second-stage stable} if admissible non-identity recoarsenings of its retained promoted surface preserve the relevant readout, signature-level closure, and control-contrast structure after the candidate is rebuilt on the coarsened surface. We say that $P$ is \emph{protocol honest} if its apparent support does not depend on pre-internalization artifacts or on reference-only protocol rows that fail the protocol gate.
\end{definition}

The joint criterion may therefore be summarized schematically as
\[
\mathrm{Struct}(P;\mathcal C)
\iff
\mathrm{Readout}(P;\mathcal C)
\wedge
\mathrm{SigUpd}(P;\mathcal C)
\wedge
\mathrm{Stable}_2(P)
\wedge
\mathrm{Honest}_{\mathrm{prot}}(P)
\wedge
\mathrm{Inv}(P)
\wedge
\mathrm{Quad}(P).
\]
This criterion is narrower than a claim of fully autonomous higher organization and broader than readout compression alone. Abstractly, the descended readout and update questions are closest to quotient and behavioral-equivalence constructions \cite{Sangiorgi2009,VanDerSchaft2004}, while the recoarsening side is closest to exact aggregation and lumpability on reduced state spaces \cite{KemenySnell1976,Buchholz1994,GivanDeanGreig2003}. The contribution is not these individual notions but the joint criterion and the positive audit results it produces on the frozen Six Birds corpus.

\begin{proposition}[Generic promoted-structure criterion]
\label{prop:generic-criterion}
Suppose a candidate promoted package $P$ on a fixed surface $\mathcal C$ satisfies the joint criterion above. Then $P$ is eligible for a package-type verdict on $\mathcal C$ as a genuinely surviving promoted object. Failure of any one of the listed conjuncts blocks that conclusion at the corresponding audit layer.
\end{proposition}

\subsection{Separate axes and nearby-support distinctions}

We keep two outcome axes separate. The first is the package-type verdict: what kind of promoted object has been earned on the frozen Route~A surface. The second is a downstream higher-layer readiness axis, recording whether the best survivor begins to participate in a weak primitive grammar.

\begin{definition}[Package verdict, readiness, and nearby-support distinction]
\label{def:outcome-axes}
The \emph{package-type verdict} is the main classification assigned by the joint internal-structure audit. The \emph{higher-layer readiness axis} records downstream weak primitive participation and is not a substitute for package type. The \emph{nearby-support distinction} records, separately, readout-surface status, detector-admissibility status, and family-level nearby characterization.
\end{definition}

\begin{proposition}[Nonreplacement boundary]
\label{prop:nonreplacement-boundary}
Evidence recorded on the higher-layer readiness axis, on the nearby-support distinction surface, or on the flagship-local chart side-branch does not by itself alter the package-type verdict. These statement surfaces answer different questions and are not collapsed into one merged label.
\end{proposition}

\begin{proposition}[Nearby-support distinctions do not collapse]
\label{prop:nearby-distinction}
Readout-surface anchor-local limitation, detector-limited nearby support, and family-level erased-partner support are logically distinct outcome surfaces in the present audit. Positive evidence on one of them does not rewrite a negative or restricted outcome on another.
\end{proposition}

The audit order and corpus are fixed in the next section.
 \section{Canonical Corpus, Candidate Classes, and Audit Pipeline}

\subsection{The frozen corpus}

All audits are performed on one frozen exact-finite corpus inherited from the previous paper rather than on a newly expanded search surface. The corpus contains two flagship anchors, four nearby positives, four required controls, and one support-only comparison row. The freeze is methodological as much as computational: we ask what internal structure is earned on the already surviving surface, not what might appear after broadening the route family or changing the comparison base.

At the level of notation, we write
\[
\mathcal C
  = \mathcal C_{\mathrm{flag}}
    \sqcup \mathcal C_{\mathrm{near}}
    \sqcup \mathcal C_{\mathrm{ctrl}}
    \sqcup \mathcal C_{\mathrm{supp}},
\qquad
\mathcal C_{\mathrm{core}}
  = \mathcal C_{\mathrm{flag}}
    \sqcup \mathcal C_{\mathrm{ctrl}}.
\]
The first decomposition records the full frozen surface used by the paper. The second isolates the core comparison subset on which the candidate classes are first placed on equal footing.

The two flagship anchors are the theorem-bearing core. The nearby positives test whether the same internal structure extends beyond the anchors on the inherited surface. The controls fix the contrast classes against which survival, honesty, and stability are read. The support-only row is available for comparison and interpretation but does not drive the package verdict. Table~\ref{tab:frozen-corpus} records the frozen corpus.

\begin{table}[t]
\centering
\caption{Frozen Route A corpus used throughout the paper.}
\label{tab:frozen-corpus}
\begin{tabularx}{\linewidth}{@{}l l >{\raggedright\arraybackslash}X@{}}
\toprule
Cohort & Case & Role in the paper \\
\midrule
Flagship & Reference-positive anchor & flagship anchor \\
Flagship & Completion-partner anchor & flagship anchor \\
Nearby positive & Nearby positive A & nearby-positive comparison case \\
Nearby positive & Nearby positive B & nearby-positive comparison case \\
Nearby positive & Nearby positive C & nearby-positive comparison case \\
Nearby positive & Nearby positive D & nearby-positive comparison case \\
Control & Null baseline & required control \\
Control & Memory-only control & required control \\
Control & Artifact/internalization control & required control \\
Control & Boundary/completion-sensitive negative & required control \\
Support-only & Erasure-recovery comparison row & support-only diagnostic \\
\bottomrule
\end{tabularx}
\end{table}

Exact internal identifiers and reproducibility hooks are carried by the artifact and reproducibility appendix; the main text refers to cases by role and human-facing label.

\subsection{Candidate classes on the frozen surface}

On that frozen corpus, we compare four candidate object types. The intrinsic candidate represents a promoted state without any additional section or protocol coordinate. The sectioned candidate represents the same promoted state together with an admissible section choice. The bundle-section-fiber candidate represents the promoted state together with an explicit admissible section fiber. The protocol-indexed candidate represents the promoted state together with explicit protocol or index data \cite{Tsiokos2026PromotionCriteria}.

These candidates are not alternative implementations of one already-assumed carrier; they are competing object hypotheses. The intrinsic candidate asks whether the surviving package is already a promoted object in its own right. The sectioned and bundle-section-fiber candidates ask whether the survivor should instead be read as charted or fibered structure. The protocol-indexed candidate asks whether support exists only after explicit internalization of protocol state. Table~\ref{tab:candidate-classes} records the candidate classes.

\begin{table}[t]
\centering
\caption{Candidate classes entering the internal-structure audit.}
\label{tab:candidate-classes}
\small
\begin{tabularx}{\linewidth}{@{}>{\raggedright\arraybackslash}p{0.18\linewidth}>{\raggedright\arraybackslash}X>{\raggedright\arraybackslash}p{0.21\linewidth}>{\raggedright\arraybackslash}X@{}}
\toprule
Candidate class & Representation carried & Additional structure & Intended role in the audit \\
\midrule
Intrinsic candidate & promoted base state & none & primary object hypothesis \\
Sectioned candidate & promoted base state + chosen admissible section & explicit section choice & charted comparison hypothesis \\
Bundle-section-fiber candidate & promoted base state + admissible section fiber & explicit section fiber & fibered comparison hypothesis \\
Protocol-indexed candidate & promoted base state + protocol/index coordinate & explicit protocol state & indexed-honesty hypothesis \\
\bottomrule
\end{tabularx}
\end{table}

The section surface inherited from the previous paper remains explicit at this stage. The intrinsic identity chart is part of the admissible comparison surface. Non-identity section families are included to motivate later comparison, but they enter only as part of the audited candidate space, not as presumptive winners. Protocol-sensitive or reference-only rows are part of the bookkeeping needed to describe the full surface but are not treated as ordinary support unless later sections show that they survive the protocol gate.

\subsection{Standardized interfaces}

To make the candidate comparison meaningful on one fixed corpus, we standardize three interfaces before presenting any substantive verdict. First, we fix a designated readout interface. The relevant readouts are unified across flagship, completion-partner, control, and support-only cases under one normalized promoted-level registry. Where the frozen surface carries exact-finite payloads, those payloads are preserved; where only source-limited records remain, the resulting rows are diagnostic only.

Second, we fix a standardized lower-generator interface for the descended update audit. At the present level of resolution, this interface is signature-level rather than a full executable promoted dynamics. The descended update question therefore asks whether a coherent promoted-level action interface is supported on the retained surface, not whether a complete promoted dynamical law has been derived.

Third, we fix an admissible recoarsening interface on the candidate's own retained surface. The relevant recoarsenings act on retained readout, section, and fiber structure and then rebuild the candidate on the coarsened surface. They are not the older vacuous operation of merging inside an already minimal promoted object. This is what allows the second-stage stability question to be asked non-vacuously.

\begin{definition}[Equal-footing comparison rule]
\label{def:equal-footing}
For a candidate kind $K$ and a case $x \in \mathcal C_{\mathrm{core}}$, the primary audit outcome is attached to the pair $(K,x)$ rather than to the raw number of admissible representatives of kind $K$ over $x$. Representative multiplicity may contribute to ambiguity statistics, but it does not count as additional positive evidence for descended readout or signature-level descended update support.
\end{definition}

Equivalently, the comparison is carried out candidate-kind by case:
\[
(K,x)
  \longmapsto
  \bigl(
    \mathbf R_K(x),\,
    \mathbf U_K(x),\,
    \mathbf S_K(x)
  \bigr),
\]
where $\mathbf R_K(x)$ records descended readout status, $\mathbf U_K(x)$ records signature-level update status, and $\mathbf S_K(x)$ records later second-stage behavior. The audit does not reward a candidate kind merely for carrying more admissible internal representatives.

\subsection{Audit order}

With the corpus and interfaces fixed, the audit proceeds in a strict order. We first compare the candidate classes under descended readout and signature-level descended update. We then subject the surviving candidates to admissible second-stage recoarsening on their own retained surfaces. We next validate invariant families and exact-finite quadratic families, separating claim-bearing families from comparison-only and proxy-only quantities. Only after that core audit do we turn to nearby-positive diagnostics, weak primitive participation, and flagship-local chart fits.

The nearby-route synthesis and downstream diagnostics are interpreted only after the main package verdict has been constrained by the core fixed-support audit. Figure~\ref{fig:audit-pipeline} summarizes this order.

\begin{figure}[t]
\centering
\includegraphics[width=0.88\linewidth]{figures/recombination/audit_pipeline.pdf}
\caption{Ordered audit logic for the paper. The main package verdict is constrained by the frozen corpus, descended-structure, second-stage, and internal-geometry audits before the nearby-route synthesis and downstream diagnostics are interpreted.}
\label{fig:audit-pipeline}
\end{figure}

We now begin that audit by comparing the candidate classes under the descended readout and signature-level update interfaces fixed here.
 \section{Descended Readout and Signature-Level Update Structure}

\subsection{Equal-footing comparison on the core subset}

The first internal-structure question is whether any candidate supports descended promoted-level readout together with a coherent promoted-level update interface on the frozen core surface. The core subset consists of the two flagship anchors and the four required controls---the comparison surface inherited from the previous paper \cite{Tsiokos2026PromotionCriteria}.

Equal footing matters because the candidate classes do not carry the same representational multiplicity. The intrinsic candidate is effectively base-level, whereas the sectioned and bundle-section-fiber candidates can carry multiple admissible representations over the same promoted state. That multiplicity contributes only to ambiguity statistics and is not counted as additional positive evidence. The relevant question is candidate kind by case: whether a candidate class supports stable descended readout and a coherent signature-level update interface on the same fixed case.

\subsection{Descended readout}

On the supported nonindexed surface, the intrinsic, sectioned, and bundle-section-fiber candidates all retain the designated promoted-level readout family on the core subset. In particular, they remain readout-sufficient on both flagship anchors and on all four required controls. At this stage, then, the candidates do not separate because some of them lose the promoted readout altogether. They separate because the same readout support has to coexist with a coherent descended action interface.

Readout sufficiency is therefore necessary but not sufficient. For sectioned and fibered candidates, the readouts must remain coherently supported across the admissible representational surface that defines the candidate, not merely carried by some representative. On the core subset, the sectioned and bundle-section-fiber candidates retain enough readout support to remain serious comparison branches. The protocol-indexed route, by contrast, does not provide supported readout on this core surface and therefore drops out of the core comparison.

\begin{proposition}[Core descended readout sufficiency]
\label{prop:descended-readout-core}
On the frozen core subset $\mathcal C_{\mathrm{core}}$, the intrinsic, sectioned, and bundle-section-fiber candidates are each readout-sufficient on all six core cases, while the protocol-indexed candidate is readout-sufficient on none of them. In particular, the nonindexed candidates remain separated at this stage not by loss of promoted readout, but by what can coexist with that readout on the same fixed cases.
\end{proposition}

\subsection{Signature-level descended update}

The update side of the audit is weaker in form but not in logical role. The retained lower-generator interface on the frozen surface is signature-level rather than a full executable promoted dynamics, so the question is whether the candidate supports a coherent descended signature-level action or closure interface, not whether a complete promoted dynamical law has been derived.

On that signature-level audit, the intrinsic candidate is the only candidate that remains closure-ready on the entire core subset. The sectioned and bundle-section-fiber candidates do not fail because the lower surface contains no usable action information; they fail because admissible section or fiber variation changes the supported operator-signature set, so the promoted update interface is not representation-stable. This is where the intrinsic candidate separates from the charted and fibered alternatives.

Note that the claim here is not a full promoted dynamics. On the frozen core surface, the intrinsic candidate alone supports both designated promoted-level readout and a coherent signature-level descended update interface; the sectioned and bundle-section-fiber candidates remain readout-sufficient but update-ambiguous.

\begin{proposition}[Core signature-level closure readiness]
\label{prop:signature-update-core}
On the same frozen core subset $\mathcal C_{\mathrm{core}}$, the intrinsic candidate is signature-level closure-ready on all six core cases. The sectioned, bundle-section-fiber, and protocol-indexed candidates are signature-level closure-ready on none of those six cases. Thus the intrinsic candidate is the only candidate alive on both descended readout and signature-level descended update on the core surface.
\end{proposition}

\subsection{Candidate comparison}

The resulting comparison is summarized in Table~\ref{tab:descended-core}. On the core subset, the intrinsic candidate is alive on both descended readout and signature-level closure readiness, while also preserving the required control contrasts. The sectioned and bundle-section-fiber candidates remain meaningful comparison branches because they preserve descended readout and control separation, but their signature-level closure interface varies across the admissible representational surface. The protocol-indexed candidate is not a supported route on this surface.

\begin{table}[t]
\centering
\caption{Core descended readout and signature-level update audit on the flagship anchors and required controls.}
\label{tab:descended-core}
\small
\begin{tabularx}{\linewidth}{@{}>{\raggedright\arraybackslash}p{0.22\linewidth}>{\raggedright\arraybackslash}p{0.17\linewidth}>{\raggedright\arraybackslash}p{0.2\linewidth}>{\raggedright\arraybackslash}p{0.14\linewidth}>{\raggedright\arraybackslash}X@{}}
\toprule
Candidate class & Core readout sufficiency & Core signature-level closure readiness & Required controls preserved & Interpretation in this section \\
\midrule
Intrinsic candidate & 6/6 core cases & 6/6 core cases & yes & alive baseline candidate \\
Sectioned candidate & 6/6 core cases & 0/6 core cases & yes & comparison branch; closure ambiguity \\
Bundle-section-fiber candidate & 6/6 core cases & 0/6 core cases & yes & comparison branch; fiber ambiguity \\
Protocol-indexed candidate & 0/6 core cases & 0/6 core cases & no & not supported on this surface \\
\bottomrule
\end{tabularx}
\end{table}

\begin{remark}[Interpretation of the first gate]
\label{rem:first-gate}
The first gate is already decisive in one restricted sense: the intrinsic route is not merely one surviving representative among several equivalent winners. It is the only candidate that clears the joint readout/update gate on the inherited core surface. The next question is whether that apparent advantage persists under admissible second-stage recoarsening, rather than disappearing once the retained promoted surface is stressed again.
\end{remark}
 \section{Second-Stage Stability and the Intrinsic Plateau}

\subsection{Admissible recoarsening on the retained surface}

The first descended-structure audit isolates the intrinsic candidate as the only class that passes both the readout and closure gates on the flagship core and required controls. That alone is not enough to certify a promoted layer: a promoted object should retain its structure when its own retained surface is recoarsened in admissible ways. The second-stage question is therefore asked on the candidate's retained promoted surface rather than by collapsing classes inside an already minimal package.

The recoarsenings used here act on retained readout, section, and fiber structure and then rebuild the candidate on the coarsened surface. This is what makes the stability test non-vacuous. The relevant question is not whether some one-chart representative can still be named after an internal merge, but whether the candidate continues to support the same core readout, closure, and control contrasts after admissible non-identity recoarsening of its own retained structure.

\subsection{The intrinsic plateau}

On that second-stage audit, the intrinsic candidate exhibits a nontrivial stability plateau: across four non-identity recoarsenings on the supported core surface, it preserves full core readout retention, full signature-level closure retention, and full required control retention. The intrinsic route is therefore not merely the first candidate to survive the initial audit; it is the only candidate whose promoted structure remains stable under admissible stress on its own retained surface.

\begin{theorem}[Intrinsic second-stage plateau]
\label{thm:intrinsic-plateau}
On the supported core surface, the intrinsic candidate has a nontrivial second-stage stability plateau. More precisely, there are four non-identity admissible recoarsenings on which the intrinsic candidate preserves full core readout retention, full signature-level closure retention, and full required-control retention.
\end{theorem}

\subsection{Fragile comparison branches and the bundle/dual check}

The sectioned and bundle-section-fiber candidates remain informative comparison branches under the same stress test, but they do not exhibit a corresponding plateau. Both retain meaningful comparison structure and continue to support readout-side information; what fails is the stable coexistence of that information with a coherent signature-level closure interface across admissible section or fiber variation. Their weakness is fragile survival rather than readout disappearance.

The same conclusion governs the bundle/dual comparison: dual augmentation does not displace the intrinsic candidate.

\begin{proposition}[Fragile comparison branches]
\label{prop:fragile-branches}
The sectioned and bundle-section-fiber candidates remain informative comparison branches under the second-stage audit, but neither carries a nontrivial stability plateau. Their surviving comparison structure is fragile rather than stable: readout-side information remains visible, while a coherent signature-level closure interface fails to persist across admissible section or fiber variation.
\end{proposition}

\begin{proposition}[Bundle/dual non-displacement]
\label{prop:bundle-dual-nondisplacement}
Dual augmentation adds no decisive improvement over the non-dual sectioned and bundle-section-fiber routes. In particular, it does not increase supported plateau size, does not repair the closure ambiguity that blocks those routes, and does not displace the intrinsic candidate from the package verdict.
\end{proposition}

\subsection{Protocol honesty and verdict boundary}

Second-stage stability alone would be insufficient if the intrinsic picture were only an artifact of omitted protocol state. The protocol gate therefore sits alongside the stability result. On the frozen surface, the intrinsic nonindexed candidate survives that gate: the core readout, closure, and stability structure is preserved on the nonindexed path, and the indexed route does not displace it. Protocol honesty does not, however, repair the nearby readout limitation; later sections return to that issue.

Taken together, the second-stage and protocol audits establish the package verdict. The intrinsic candidate is the only route that passes the first descended-structure gate, remains stable under nontrivial second-stage recoarsening, and survives protocol treatment without being displaced by sectioned, bundle, dual, or indexed alternatives. Figure~\ref{fig:second-stage-stability} and Table~\ref{tab:stability-protocol} summarize this contrast.

\begin{proposition}[Protocol-honest nonindexed preservation]
\label{prop:protocol-honest-intrinsic}
The intrinsic nonindexed route passes the protocol gate on the frozen surface. The indexed route does not displace it, and protocol treatment does not convert the nearby readout limitation into a repaired neighborhood statement.
\end{proposition}

\begin{remark}[Verdict boundary]
\label{rem:stability-boundary}
This section carries the package verdict most directly; invariant and quadratic content are addressed in the next section, and the nearby readout limitation is carried forward in its own diagnostic layer.
\end{remark}

\begin{table}[t]
\centering
\small
\begin{tabularx}{\linewidth}{@{}>{\raggedright\arraybackslash}p{0.20\linewidth}>{\raggedright\arraybackslash}p{0.14\linewidth}c>{\raggedright\arraybackslash}X>{\raggedright\arraybackslash}p{0.16\linewidth}@{}}
\toprule
Candidate class & Second-stage status & Plateau size & Protocol status & Role in the verdict \\
\midrule
Intrinsic candidate & stable plateau & 4 & honest nonindexed support preserved & main verdict candidate \\
Sectioned candidate & fragile survival & 0 & comparison-only / no displacement & comparison branch \\
Bundle-section-fiber candidate & fragile survival & 0 & comparison-only / no displacement & comparison branch \\
Protocol-indexed candidate & not decisive & 0 & no displacement of nonindexed intrinsic route & indexed comparison route \\
\bottomrule
\end{tabularx}
\caption{Second-stage stability and protocol-honesty comparison on the promoted candidate classes.}
\label{tab:stability-protocol}
\end{table}

\begin{figure}[t]
\centering
\includegraphics[width=0.82\linewidth]{figures/recombination/second_stage_stability.pdf}
\caption{Second-stage stability contrast across promoted candidate classes. Only the intrinsic route has a nontrivial supported plateau; sectioned and bundle/fiber routes remain fragile comparison branches, and the indexed route does not displace the honest nonindexed intrinsic core.}
\label{fig:second-stage-stability}
\end{figure}
 \section{Retained Invariants and Exact-Finite Quadratic Structure}

\subsection{Object-map validation and retained invariant families}

Second-stage stability identifies the intrinsic plateau as the retained promoted surface on which internal geometry should be sought. The next question is what, exactly, is retained there. In this paper, that question is not answered by simply listing quantities with good numerical behavior. Any candidate invariant must first survive an object-map audit: it must be clear what retained promoted-level behavior the quantity is tracking, and proxy bookkeeping must be kept distinct from the object itself.

On that audit, four invariant families remain claim-bearing on the intrinsic plateau. The first is a plateau-retention signature, which records the non-identity recoarsenings under which the intrinsic candidate preserves full core readout, signature-level closure, and required control behavior. The second is a core readout-support signature, which tracks stable promoted readout support on the retained intrinsic surface. The third is a core operator-support signature, which tracks stable signature-level descended action support on that same surface. The fourth is a control-contrast signature, which records the continued separation of the required controls from the flagship-positive intrinsic surface.

These families are retained because they track retained promoted-level structure rather than merely counting representatives or recording accidental complexity. They are exact or composite-exact on the frozen surface and remain anchored to the same intrinsic plateau that carried the package verdict in the previous section. Their role is therefore not merely descriptive: together they identify the stable promoted surface on which later geometric probes can be interpreted as internal structure rather than as one-shot numerical coincidences.

For compact reference, we write
\[
\mathcal I_{\mathrm{ret}}
  =
  \{
    \text{plateau-retention},
    \text{core readout-support},
    \text{core operator-support},
    \text{control-contrast}
  \}.
\]

\begin{proposition}[Retained invariant families]
\label{prop:retained-invariants}
On the intrinsic plateau, the four invariant families in $\mathcal I_{\mathrm{ret}}$ are retained as claim-bearing invariant content. Each tracks a retained promoted-level feature of the stable intrinsic surface, and none reduces to raw representative counting or proxy bookkeeping.
\end{proposition}

\subsection{Comparison-only and proxy-only families}

Not every useful family is claim-bearing. The section/fiber ambiguity signature remains exact and informative, but its role is diagnostic: it tracks how section or fiber variation changes the supported readout or operator interface, helping explain why those branches remain fragile under stress, without converting them into stable promoted winners.

The same discipline applies to proxy-only quantities. Bundle/fiber cardinality or bundle-complexity profiles may summarize representational burden, but they are not the bundle object itself. Count-based descriptors remain useful as bookkeeping, yet they are not promoted to invariant status. Internal geometry is claimed only where the object map has been validated and the retained quantity tracks a stable promoted-level feature rather than a numerical shadow of representation size.

\subsection{Exact-finite quadratic families}

On the intrinsic plateau, the strongest positive internal-geometry signal comes from exact-finite quadratic structure. Three quadratic families are retained as claim-bearing. The first is the self-quadratic readout family, which records the intrinsic exact-finite readout mass in a positive quadratic form. The second is a gap-type quadratic family that records exact-finite readout separation through squared pairwise differences. The third is a recoarsening-consistency quadratic family that measures exact-finite compatibility between the intrinsic retained surface and its nontrivial plateau recoarsenings.

The positive quadratic families isolated here sit closest to exact finite energy constructions---graph-Laplacian quadratic forms, Dirichlet forms, and carr\'{e}-du-champ-style objects---although their role here is internal to the promoted-package audit rather than imported from a separate analytic theory \cite{Chung1997,FukushimaOshimaTakeda1994,BakryEmery1985}.

These retained quadratic families share the same structural profile. They are exact-finite rather than merely heuristic; they are nonnegative on the evaluated intrinsic plateau surface; they remain honest under the protocol gate; they are stable on the intrinsic nontrivial plateau rather than only on one favorable recoarsening; and they continue to separate the required controls from the flagship-positive intrinsic surface. For the purposes of this paper, that is the sense in which exact-finite quadratic structure supplies a real internal-geometry signal.

By contrast, not every quadratic probe is promoted to that role. Section/fiber variation quadratics remain valuable diagnostics for the comparison branches, but they are not the paper's main positive geometric object. Likewise, unresolved control-distance or other non-retained probe families are recorded as unresolved or non-claim-bearing and are not used to extend the main verdict.

We summarize the claim-bearing quadratic surface as
\[
\mathcal Q_{\mathrm{ret}}
  =
  \{
    q_{\mathrm{readout\_l2\_self}},
    q_{\mathrm{readout\_gap}},
    q_{\mathrm{coarsening\_consistency}}
  \}.
\]

\begin{proposition}[Retained exact-finite quadratic families]
\label{prop:retained-quadratics}
On the intrinsic plateau, the three families in $\mathcal Q_{\mathrm{ret}}$ are retained as claim-bearing exact-finite quadratic structure. They are exact-finite, remain honest under the protocol gate, persist across the nontrivial intrinsic plateau, and do not reduce to comparison-only ambiguity probes or unresolved distance summaries.
\end{proposition}

\subsection{Interpretation of the internal-geometry signal}

On the frozen exact-finite intrinsic plateau, there exists a retained family of validated invariant content together with a retained family of exact-finite quadratic probes that behave as promoted-level internal structure. This does not embed the intrinsic plateau in any ambient carrier, nor does it promote every useful comparison statistic to invariant status. Tables~\ref{tab:invariant-families} and \ref{tab:quadratic-families} summarize which families are claim-bearing, comparison-only, proxy-only, or unresolved on the retained surface.

\begin{remark}[What this section does not claim]
\label{rem:no-global-linearity}
This section does not claim any ambient additive or linear carrier model, does not promote bundle/dual material to a competing geometry verdict, and does not convert comparison-only or proxy-only families into claim-bearing structure.
\end{remark}

\begin{table}[t]
\centering
\small
\begin{tabularx}{\linewidth}{@{}>{\raggedright\arraybackslash}p{0.22\linewidth}>{\raggedright\arraybackslash}p{0.12\linewidth}>{\raggedright\arraybackslash}X>{\raggedright\arraybackslash}p{0.22\linewidth}@{}}
\toprule
Family & Status & Target object & Role in the paper \\
\midrule
Plateau-retention signature & retained & plateau membership / retained promoted surface & claim-bearing intrinsic invariant \\
Core readout-support signature & retained & descended readout support & claim-bearing intrinsic invariant \\
Core operator-support signature & retained & signature-level closure/readout interface support & claim-bearing intrinsic invariant \\
Control-contrast signature & retained & control separation on the retained surface & claim-bearing intrinsic invariant \\
Section/fiber ambiguity signature & comparison-only & representational ambiguity across section/fiber branches & comparison diagnostic \\
Bundle/fiber cardinality profile & proxy-only & representation-size bookkeeping only & non-claim-bearing \\
\bottomrule
\end{tabularx}
\caption{Validated invariant families on the retained promoted surface.}
\label{tab:invariant-families}
\end{table}

\begin{table}[t]
\centering
\small
\begin{tabularx}{\linewidth}{@{}>{\raggedright\arraybackslash}p{0.26\linewidth}>{\raggedright\arraybackslash}p{0.10\linewidth}>{\raggedright\arraybackslash}X>{\raggedright\arraybackslash}p{0.18\linewidth}@{}}
\toprule
Family & Status & Structural profile & Role in the paper \\
\midrule
Self readout quadratic (\texttt{q\_readout\_l2\_self}) & retained & exact-finite, nonnegative, plateau-stable, protocol-honest & claim-bearing intrinsic geometry \\
Gap quadratic (\texttt{q\_readout\_gap}) & retained & exact-finite, nonnegative, control-separating & claim-bearing intrinsic geometry \\
Recoarsening-consistency quadratic (\texttt{q\_coarsening\_consistency}) & retained & exact-finite, plateau-stable, control-retaining & claim-bearing intrinsic geometry \\
Section/fiber variation quadratic & comparison-only & exact diagnostic of branch ambiguity & comparison diagnostic \\
Control-distance quadratic & unresolved & not promoted on the frozen surface & non-claim-bearing \\
\bottomrule
\end{tabularx}
\caption{Exact-finite quadratic families and their status in the paper.}
\label{tab:quadratic-families}
\end{table}
 \section{Nearby Positives: Detector-Limited Readout and Erased-Partner Family Support}

\subsection{The readout-surface limitation}

On the frozen exact readout surface, the intrinsic route remains anchor-only: the nearby positives do not inherit the flagship exact readout support. Two questions must therefore be distinguished. The first is whether the current exact readout lens extends to the nearby positives; on the frozen surface the answer is no. The second is whether the nearby positives are thereby outside the same promoted family altogether. The diagnostics assembled here address the second question without erasing the first.

\begin{proposition}[Readout-surface anchor-only limitation]
\label{prop:nearby-anchor-only}
On the frozen exact and decision-bearing readout surface, the nearby positives remain anchor-only. The current exact readout lens does not extend from the flagship anchors to the nearby-positive cases.
\end{proposition}

\subsection{Carrier support and detector admissibility}

The first nearby diagnostic is detector-independent: a carrier-level closure suite calibrated from the required controls.

\begin{proposition}[Carrier support present]
\label{prop:nearby-carrier-support}
On the carrier-level closure suite calibrated from the required controls, all four nearby positives exhibit positive carrier support rather than carrier emptiness. Hence the nearby cases are not negative in the same sense as the control surface.
\end{proposition}

The nearby positives remain structurally related to the surviving flagship object even though the current exact readout lens does not certify them directly. The next question is whether a different exact-finite detector could unlock them on the same frozen readout surface. The restriction against detector drift is methodological rather than cosmetic: it follows the same logic as standard concerns about undisclosed analytic flexibility and valid post-selection inference \cite{SimmonsNelsonSimonsohn2011,GelmanLoken2013,BerkBrownBujaZhangZhao2013}.

\begin{proposition}[Detector unlock remains restricted]
\label{prop:nearby-detector-restricted}
No additional theorem-admissible exact-finite nearby detector unlock exists on the frozen readout surface. The nearby failure mode on that surface is detector-limited and source-limited rather than an absence of all nearby structure.
\end{proposition}

Accordingly, the anchor-only outcome is a limitation of the present exact readout interface, not a finding of nearby object absence.

\subsection{Family-level quotient support}

The readout-surface question is not the only family-level question available on fixed support. The family-level quotient strictly refines the native or branchwise quotient, and the nearby cases can be audited against it on admissible exact family-level witness evidence. The marked/erased vocabulary is borrowed from the erasure tradition, but is used here only to label a family-level quotient pattern on fixed support rather than to import that formalism wholesale \cite{ScullyDruhl1982}.

\begin{proposition}[Erased-partner family support]
\label{prop:nearby-erased-partner}
At the family-level quotient layer, all four nearby positives support an erased-partner interpretation rather than a negative one. This classification is not obtained by post hoc detector replacement, and it is not explained by protocol or completion artifacts.
\end{proposition}

This classification is casewise. It shows that all four nearby positives match the erased-partner pattern on the frozen surface, but it does not by itself identify them as one and the same erased object or as partners of one single submechanism.

\subsection{No remaining unresolved nearby cases}

Once the family-level quotient audit is taken into account, no unresolved nearby cases remain on the frozen admissible surface, and the candidate-subset witness search is not needed. The remaining open point is what the current exact readout surface can and cannot certify directly.

\begin{remark}[Subset witness search not needed]
\label{rem:no-nearby-witness-search}
The candidate-subset witness search is skipped on the current frozen nearby surface because no unresolved nearby-positive cases remain after the family-level erased-partner classification is taken into account.
\end{remark}

\subsection{Synthesis and caveat}

The nearby-positive synthesis is layered: anchor-only on the frozen exact readout surface, structurally positive on detector-independent carrier diagnostics, erased-partner on the family-level quotient audit, and detector-restricted on the admissibility audit. These statements answer different structural questions and are not contradictory. The nearby route closes with erased-partner support rather than with object absence, under the caveat that readout-level detector unlock remains restricted. Figure~\ref{fig:nearby-synthesis} and Tables~\ref{tab:nearby-route-synthesis} and \ref{tab:nearby-positive-classification} summarize the layered nearby route.

\begin{proposition}[Nearby-route synthesis with caveat]
\label{prop:nearby-caveat}
The nearby route closes positively but with caveat. On the frozen exact readout surface the nearby positives remain anchor-only; on detector-independent carrier diagnostics they remain positive; at the family-level quotient layer they support erased-partner interpretation; and no additional theorem-admissible exact-finite detector unlock exists on the current frozen readout surface. Hence the nearby route closes with preserved caveat: readout-level detector unlock remains restricted while family-level erased-partner support is positive.
\end{proposition}

\begin{table}[t]
\centering
\small
\begin{tabularx}{\linewidth}{@{}>{\raggedright\arraybackslash}p{0.19\linewidth}>{\raggedright\arraybackslash}p{0.13\linewidth}>{\raggedright\arraybackslash}X>{\raggedright\arraybackslash}p{0.22\linewidth}@{}}
\toprule
Audit layer & Result & What it establishes & What it does not establish \\
\midrule
Readout-surface test & anchor-only & current exact readout lens does not extend to nearby positives & does not show object absence \\
Carrier-level closure suite & carrier support present & nearby positives are not carrier-empty on control-frozen carrier metrics & does not unlock a new theorem-admissible detector \\
Detector admissibility audit & detector unlock restricted & no additional theorem-admissible exact-finite nearby unlock on the frozen readout surface & does not imply nearby negatives \\
Family-level quotient audit & erased-partner support & nearby positives match the erased-partner pattern on admissible quotient evidence & does not rewrite the readout-surface limitation \\
Subset witness search & not needed & all nearby cases resolved at quotient layer & does not add a separate local witness claim \\
Final synthesis & proceed with caveat & nearby route closes positively with erased-partner support & does not alter the package verdict \\
\bottomrule
\end{tabularx}
\caption{Nearby-positive route synthesis on the frozen admissible surface.}
\label{tab:nearby-route-synthesis}
\end{table}

\begin{table}[t]
\centering
\begin{tabular}{p{0.18\textwidth}p{0.15\textwidth}p{0.15\textwidth}p{0.18\textwidth}p{0.24\textwidth}}
\hline
Nearby case & Readout-surface status & Carrier-level status & Family-level quotient status & Overall interpretation in this section \\
\hline
Nearby positive A & anchor-only & carrier support & erased-like & detector-limited erased-partner support \\
Nearby positive B & anchor-only & carrier support & erased-like & detector-limited erased-partner support \\
Nearby positive C & anchor-only & carrier support & erased-like & detector-limited erased-partner support \\
Nearby positive D & anchor-only & carrier support & erased-like & detector-limited erased-partner support \\
\hline
\end{tabular}
\caption{Nearby-positive classification on the frozen admissible surface.}
\label{tab:nearby-positive-classification}
\end{table}

\begin{figure}[t]
\centering
\includegraphics[width=0.84\linewidth]{figures/recombination/nearby_synthesis.pdf}
\caption{Layered nearby-route synthesis on the frozen admissible surface. The readout-surface anchor-only limitation is retained, while carrier and family-level diagnostics support a positive erased-partner interpretation that closes only with caveat.}
\label{fig:nearby-synthesis}
\end{figure}
 \section{Downstream Diagnostics: Weak Primitive Participation and Flagship-Local Charts}

\subsection{Weak primitive participation}

With the package-type verdict fixed, the remaining question is whether the best surviving candidate begins to participate, even weakly, in a higher-layer primitive grammar. The weak primitive-descent stress panel applied here is downstream: it asks what additional action-level support the surviving package may already exhibit, not what kind of promoted package has survived.

On that panel, the intrinsic candidate is the only route with positive support across the core promoted-action slots. The sectioned and bundle-section-fiber candidates remain comparison branches. Exploratory rewrite/extension or budget-style slots are recorded as annex material only.

\begin{proposition}[Weak primitive participation and readiness]
\label{prop:readiness}
On the weak primitive-descent core panel, the intrinsic candidate is the only surviving route with positive support across all core primitive slots. This supports a separate higher-layer readiness result, recorded on a downstream axis, and does not replace the package-type verdict.
\end{proposition}

\subsection{Flagship-local chart fits}

On the two-flagship patch, the intrinsic candidate admits exact finite local affine, additive, and matrix-like fits that match the designated readout data and the signature-level action interface. Local realization and linearization ideas are classical \cite{Sontag1979,GuckenheimerHolmes1983}; the point here is only that such a fit appears on the retained flagship patch in exact finite form.

\begin{proposition}[Flagship-local chart support]
\label{prop:local-chart}
On the flagship patch, the intrinsic candidate admits exact local affine/additive/matrix-like fits that match both the designated readout data and the signature-level action interface. This is a flagship-local, non-gating compatibility result; it does not extend to the nearby positives on the same exact-finite surface and does not imply any ambient additive or linear carrier model.
\end{proposition}

\subsection{Summary}

Weak primitive participation and a flagship-local chart fit enrich the picture established by the core internal-structure audit but do not alter the package verdict. Tables~\ref{tab:primitive-readiness} and \ref{tab:local-chart} summarize these downstream results.

\begin{remark}[Downstream nonreplacement]
\label{rem:downstream-nonreplacement}
Neither higher-layer readiness nor flagship-local chart support changes the package verdict, repairs the nearby readout limitation, or licenses any claim of ambient linear structure or a fully autonomous higher-level layer.
\end{remark}

\begin{table}[t]
\centering
\small
\begin{tabularx}{\linewidth}{@{}>{\raggedright\arraybackslash}p{0.20\linewidth}>{\raggedright\arraybackslash}X>{\raggedright\arraybackslash}p{0.16\linewidth}>{\raggedright\arraybackslash}p{0.22\linewidth}@{}}
\toprule
Candidate class & Core weak primitive result & Evidence level & Higher-layer role \\
\midrule
Intrinsic candidate & weak passes on all core primitive slots & genuinely promoted-level & supports separate level-2 readiness axis \\
Sectioned candidate & no core weak pass & comparison-only & comparison branch only \\
Bundle-section-fiber candidate & no core weak pass & comparison-only & comparison branch only \\
Protocol-indexed candidate & not a decisive surviving route & not claim-bearing & not used to set readiness \\
\bottomrule
\end{tabularx}
\caption{Downstream primitive-descent and higher-layer readiness summary.}
\label{tab:primitive-readiness}
\end{table}

\begin{table}[t]
\centering
\small
\begin{tabularx}{\linewidth}{@{}>{\raggedright\arraybackslash}p{0.20\linewidth}>{\raggedright\arraybackslash}X c c >{\raggedright\arraybackslash}p{0.22\linewidth}@{}}
\toprule
Candidate class & Chart status & Readout fit & Signature fit & Role in the paper \\
\midrule
Intrinsic candidate & exact local chart present & exact & exact & supportive, non-gating flagship-local result \\
Sectioned candidate & comparison-only & not decisive & not decisive & comparison-only side-branch \\
Bundle-section-fiber candidate & partial comparison fit & partial & partial & comparison-only side-branch \\
\bottomrule
\end{tabularx}
\caption{Local-chart side-branch on the flagship patch.}
\label{tab:local-chart}
\end{table}
 \section{Formal Criterion Layer}

\subsection{Current formalized vocabulary}

The computational audits are accompanied by a Lean formal layer whose purpose is to stabilize the vocabulary in which the internal-structure claims are stated, not to re-prove the empirical atlas. The current layer is implemented in Lean~4 with support from mathlib \cite{deMouraUllrich2021,MathlibCommunity2020}, but its role here is criterion-level rather than theorem-exhaustive. The layer formalizes the abstract types and interfaces for candidate promoted packages, admissible section structure, descended readout, descended update, recoarsening, and second-stage stability.

The formal surface distinguishes the four candidate kinds used in the computational audit, contains abstract package and section-family vocabulary, and provides descended readout, descended update, and second-stage stability interfaces phrased at the level of abstract compatibility rather than as a concrete replay of the computational experiment.

\begin{definition}[Current formal criterion layer]
\label{def:formal-layer}
The \emph{current formal criterion layer} is the Lean surface consisting of abstract vocabulary for candidate kinds, candidate packages, section families, descended readout, descended update, recoarsening, and second-stage stability, together with criterion-level statement surfaces for package verdict, higher-layer readiness, nearby-support distinctions, and nonreplacement boundaries.
\end{definition}

\subsection{Criterion layer and separate axes}

On top of that definitional layer, the formal surface adds criterion vocabulary: package-type verdict, higher-layer readiness, intrinsic versus protocol-indexed outcomes, bundle/dual upgrade, and nearby-support distinctions. In the formal layer, readout-surface anchor-local limitation, detector-limited nearby support, erased-partner family-level support, and nonreplacement boundaries are represented as distinct statement-level objects, mirroring the separation of outcome axes in the audit design.

\begin{proposition}[Current formal scope]
\label{prop:formal-scope}
The current Lean layer formalizes the abstract criterion vocabulary and the separation of outcome axes used in the manuscript, but it does not prove the empirical package verdict, the nearby-route synthesis, the flagship-local chart result, or the downstream readiness result.
\end{proposition}

\subsection{Flagship stubs and current scope}

The formal layer also contains dedicated flagship stubs that package the abstract ingredients needed for later flagship-level theorem work---a candidate-package surface, the relevant section structure, descended readout and update interfaces, and stability-oriented criterion targets---without providing empirical witness theorems for the reference-positive or completion-partner anchors.

\subsection{Proof-gap boundary}

What is presently formalized: abstract definitions, criterion surfaces, separation of outcome axes, and flagship stub interfaces. What remains open: substantive theorems connecting those abstractions to concrete empirical verdicts, and any stronger formal treatment of nearby-support synthesis, protocol-indexed versus nonindexed classification, bundle/dual upgrade claims, or local-chart consequences. The formal layer does not assert as a theorem that the package verdict is intrinsic or that the nearby route closes positively; it records the vocabulary in which such claims could later be stated and proved. Table~\ref{tab:formal-layer} summarizes the current scope.

\begin{remark}[Proof-gap boundary]
\label{rem:formal-proof-gap}
The formal layer supplies definitions, criteria, flagship stub interfaces, and proof-gap markers, but does not turn the empirical atlas into a proved Lean theorem.
\end{remark}

\begin{table}[t]
\centering
\small
\begin{tabularx}{\linewidth}{@{}>{\raggedright\arraybackslash}p{0.14\linewidth}>{\raggedright\arraybackslash}X>{\raggedright\arraybackslash}p{0.22\linewidth}>{\raggedright\arraybackslash}p{0.22\linewidth}@{}}
\toprule
Component & Current formal content & Role in the paper & Not yet claimed \\
\midrule
Definitions & candidate-package, section-family, descended readout/update, recoarsening/stability vocabulary & stabilizes the paper's abstract language & no empirical verdicts proved \\
Criteria & package verdict axis, higher-layer readiness axis, intrinsic/indexed and bundle/dual criterion surfaces & keeps outcome axes separate & no final classification theorem proved \\
Flagship stubs & abstract flagship witness interfaces & prepares later flagship-level theorem work & no concrete flagship witness theorem proved \\
Proof-gap layer & explicit open theorem targets and caveat boundaries & localizes what remains future work & no nearby-route or local-chart theorem claimed \\
\bottomrule
\end{tabularx}
\caption{Current Lean formal layer and its role in the paper.}
\label{tab:formal-layer}
\end{table}
 \section{Discussion and Limitations}

\subsection{Why the verdict is intrinsic}

The intrinsic nonindexed candidate is the only candidate that survives the entire claim-bearing audit stack: the joint descended readout and signature-level update gate (\Cref{prop:descended-readout-core,prop:signature-update-core}), the nontrivial second-stage plateau (\Cref{thm:intrinsic-plateau}), and the retained invariant and quadratic families (\Cref{prop:retained-invariants,prop:retained-quadratics}). The sectioned and bundle-section-fiber candidates remain informative---their failure mode is closure ambiguity under section or fiber variation, not readout collapse---but they do not supply a stable promoted object. Dual augmentation adds no decisive improvement, and the protocol-indexed route does not displace the nonindexed intrinsic core. The verdict is the result of elimination on fixed support, not a choice among equally good summaries.

\subsection{Nearby route and caveat}

The nearby positives do not inherit the flagship exact readout support on the frozen surface, but they are not thereby reduced to negatives. Carrier diagnostics show they are not carrier-empty, and the family-level quotient audit supports an erased-partner interpretation (\Cref{prop:nearby-anchor-only,prop:nearby-carrier-support,prop:nearby-detector-restricted,prop:nearby-erased-partner,prop:nearby-caveat}). The readout-level detector unlock remains restricted while family-level erased-partner support is positive. This prevents the nearby route from being read either as a clean neighborhood-coherence theorem or as a failed family extension.

\subsection{Downstream diagnostics}

The primitive-descent stress panel and the flagship-local chart fit (\Cref{prop:readiness,prop:local-chart}) enrich the picture but do not alter the package verdict. Weak higher-layer readiness is not a second package verdict, and a flagship-local chart is not an ambient-carrier theorem.

\subsection{Limitations}

Several limitations apply. First, the update surface remains signature-level rather than a full executable promoted dynamics. Second, the nearby-positive readout gap remains real on the current exact frozen readout surface; the erased-partner synthesis explains that gap but does not remove it. Third, no ambient additive or linear carrier model is claimed, and the flagship-local chart is not treated as evidence of such a carrier. Fourth, the bundle/dual and indexed routes remain comparison branches rather than final verdicts. Fifth, the formal layer is a criterion and statement surface rather than a full mechanization of the empirical atlas.

Together these limitations define the boundary of the present claim: a stable intrinsic promoted package with internal structure on the frozen exact-finite surface, a nearby route that closes positively but with caveat, and downstream diagnostics that remain non-gating.
 \section{Conclusion}

On the frozen exact-finite corpus, the surviving promoted object on the fixed-support Route~A flagship surface is an intrinsic package. The intrinsic nonindexed candidate is the only candidate that combines descended readout, signature-level descended update support, nontrivial second-stage stability, protocol honesty, retained invariant structure, and retained exact-finite quadratic structure without being displaced by sectioned, bundle, dual, or indexed alternatives (\Cref{prop:signature-update-core,thm:intrinsic-plateau,prop:protocol-honest-intrinsic,prop:retained-invariants,prop:retained-quadratics}).

The nearby-positive route closes positively but with caveat. On the frozen exact readout surface, the nearby positives remain anchor-only; carrier-level and family-level diagnostics show that they support an erased-partner interpretation rather than a negative one, under the restriction that no further theorem-admissible exact-finite detector unlock is available on the current frozen surface (\Cref{prop:nearby-caveat}).

The Six Birds substrate supports an intrinsic promoted package with stable internal structure on the Route~A flagship family under fixed support. The layered structure built by the series is therefore not merely a bookkeeping convenience: the substrate produces a genuine higher-level object with its own descended organization. Weak higher-layer readiness and a flagship-local exact chart are downstream diagnostics. The nearby extension and local representability of the survivor remain bounded by the frozen readout surface.
 
\appendix
\section{Formal Definitions and Theorem Statements}
\label{app:formal-definitions}

This appendix collects the abstract formal vocabulary used by the paper without turning the empirical audit into a theorem by declaration. The goal is to make the statement surface explicit and conventional and to record, in one place, the definition and proposition surfaces used in the main text.

\subsection*{Abstract definitions}

The main text introduces the following definition surfaces.
\begin{itemize}
\item \Cref{def:candidate-package}: candidate promoted package.
\item \Cref{def:descended-interfaces}: descended readout support and signature-level descended update support.
\item \Cref{def:stability-honesty}: second-stage stability and protocol honesty.
\item \Cref{def:outcome-axes}: package verdict, higher-layer readiness, and nearby-support distinction.
\item \Cref{def:equal-footing}: equal-footing comparison rule on the core subset.
\item \Cref{def:formal-layer}: the current formal criterion layer used by the bounded Lean track.
\end{itemize}

\subsection*{Criterion surfaces}

On top of those definitions, the Lean layer records criterion names rather than empirical verdict theorems. The main criterion objects are:
\begin{itemize}
\item \texttt{PackageTypeVerdict}, \texttt{HigherLayerReadiness}, \texttt{NearbySupportCharacterization}, and \texttt{ReadoutSurfaceNeighborhoodStatus} as separate outcome axes.
\item \texttt{PromotedStructureSuccess}, \texttt{HonestIntrinsicCriterion}, \texttt{ProtocolIndexedCriterion}, and \texttt{BundleDualUpgradeCriterion} as abstract criterion surfaces.
\item \texttt{ReadoutSurfaceAnchorOnly}, \texttt{DetectorLimitedNearbySupport}, \texttt{ErasedPartnerFamilySupport}, and \texttt{RespectsNonreplacementBoundary} as statement-level distinction surfaces.
\end{itemize}

\subsection*{Main statement surface}

The main text packages the empirical claim surfaces by propositions, a theorem, and remarks rather than by prose alone. The load-bearing statement surfaces are:
\begin{itemize}
\item \Cref{prop:generic-criterion,prop:nonreplacement-boundary,prop:nearby-distinction} for the criterion layer.
\item \Cref{prop:descended-readout-core,prop:signature-update-core} for the first descended-structure gate.
\item \Cref{thm:intrinsic-plateau,prop:fragile-branches,prop:bundle-dual-nondisplacement,prop:protocol-honest-intrinsic} for second-stage stability and protocol honesty.
\item \Cref{prop:retained-invariants,prop:retained-quadratics} for retained internal geometry.
\item \Cref{prop:nearby-anchor-only,prop:nearby-carrier-support,prop:nearby-detector-restricted,prop:nearby-erased-partner,prop:nearby-caveat} for the nearby-route closeout.
\item \Cref{prop:readiness,prop:local-chart,rem:downstream-nonreplacement} for downstream diagnostics.
\item \Cref{prop:formal-scope,rem:formal-proof-gap} for the bounded formal layer.
\end{itemize}

\subsection*{Lean boundary}

What is currently proved in Lean is light: projection lemmas, definitional identities, helper facts about descended interfaces, finite-type separations, and a constructor-style lemma for the second-stage stability predicate. No empirical theorem is claimed for the intrinsic verdict, the nearby-route synthesis, the higher-layer readiness result, or the flagship-local chart side-branch.
 \section{Computational Methods}
\label{app:computational-methods}

This appendix summarizes the computational method discipline behind the frozen audit surface used throughout the paper. It is included for auditability, not to enlarge the paper's claim boundary.

\subsection*{Frozen corpus and candidate interfaces}

All decision-bearing computations are carried out on one frozen corpus
\[
\mathcal C
  = \mathcal C_{\mathrm{flag}}
    \sqcup \mathcal C_{\mathrm{near}}
    \sqcup \mathcal C_{\mathrm{ctrl}}
    \sqcup \mathcal C_{\mathrm{supp}},
\qquad
\mathcal C_{\mathrm{core}}
  = \mathcal C_{\mathrm{flag}}
    \sqcup \mathcal C_{\mathrm{ctrl}},
\]
with two flagship anchors, four nearby positives, four required controls, and one support-only comparison row. Candidate interfaces are standardized before verdict-level comparison. The live candidate classes are the intrinsic, sectioned, bundle-section-fiber, and protocol-indexed routes described in the main text.

\subsection*{Equal-footing comparison rule}

Candidate comparison is performed on an equal-footing candidate-kind-by-case aggregation rather than on raw instance multiplicity. In the notation of \Cref{def:equal-footing}, the primary audit attaches to $(K,x)$ with $K$ a candidate kind and $x \in \mathcal C_{\mathrm{core}}$. This is why section or fiber multiplicity can contribute to ambiguity diagnostics without counting as extra positive evidence. Source-limited and reference-only rows are kept explicit and do not silently become decision-bearing exact support.

\subsection*{Readout, update, and recoarsening interfaces}

The readout interface is exact-finite where the frozen surface supports exact payloads and diagnostic-only where only source-limited support remains. The update surface is signature-level rather than a full executable promoted dynamics. The recoarsening interface acts on the retained promoted surface and then rebuilds the candidate on the coarsened surface; it does not reopen the vacuous merge-inside-minimal-package route excluded by the second-stage audit. This is the manuscript-level implementation of the criteria recorded in \Cref{def:descended-interfaces,def:stability-honesty}.

\subsection*{Invariant validation and nearby-route methodology}

Invariant and quadratic families are screened through object-map validation and proxy discipline before they become claim-bearing. Comparison-only and proxy-only families remain visible but demoted. The nearby route is audited separately through readout-surface support, carrier closure, detector admissibility, family-level quotient structure, and the conditional subset-witness fallback. This is why the nearby-positive closeout can preserve the readout-surface anchor-only limitation while still reporting positive erased-partner support under caveat.
 \section{Artifact and Reproducibility Note}
\label{app:repro}

The frozen paper surface is indexed and summarized in the repo's reproducibility materials. The purpose of this appendix is to give a concise rerun and verification path for the auditable frozen surface used by the paper.

\subsection*{Rebuild and verification commands}

\begin{verbatim}
./scripts/build_paper.sh
cd lean && lake build
\end{verbatim}

The repository reproducibility manifest records the frozen-pipeline execute and
verify commands for the artifact surface summarized here. The compact command
sequence above is the manuscript-facing rebuild path used for this paper.

\subsection*{Final frozen state}

\begin{itemize}
\item package-type verdict: \texttt{honest\_intrinsic\_package}
\item higher-layer readiness: \texttt{level2\_supported}
\item best surviving candidate: \texttt{intrinsic\_naked / nonindexed}
\end{itemize}

\subsection*{Nearby-route closeout}

The nearby-route closeout remains layered rather than collapsed:
\begin{itemize}
\item readout-surface status: \texttt{anchor\_only}
\item route closeout: \texttt{nearby\_route\_exhausted\_positive}
\item freeze recommendation: \texttt{proceed\_with\_caveat}
\item caveat: readout-level detector unlock remains restricted while family-level erased-partner support is positive
\end{itemize}

The nearby-route diagnostic chain is part of the final auditable surface and is carried forward as supporting diagnostics rather than as a replacement for the package verdict.

\subsection*{Frozen surface categories}

The frozen asset index records the following categories for the paper-facing surface:
\begin{itemize}
\item \texttt{final\_core}: 2
\item \texttt{final\_support}: 26
\item \texttt{final\_diagnostic}: 6
\item \texttt{manifest\_or\_note}: 14
\item \texttt{optional\_conditional}: 2
\end{itemize}
 \section{Lean Formalization Summary}
\label{app:lean-summary}

The Lean surface imported through \texttt{Recombination.Promotion} is compile-clean. It provides a stable formal vocabulary for the paper's criterion layer; the exact-finite empirical atlas has not been mechanized.

\subsection*{File-level summary}

\begin{itemize}
\item \texttt{Definitions.lean} contains the abstract package, section, readout, update, recoarsening, and outcome-axis vocabulary.
\item \texttt{Criteria.lean} adds the criterion predicates and the separate nearby-support and nonreplacement surfaces.
\item \texttt{Statements.lean} contains only trivial helper lemmas and projection facts.
\item \texttt{FlagshipStub.lean} packages abstract flagship-oriented interfaces for later theorem wiring.
\item \texttt{ProofGaps.lean} records the proved surface, stub surface, open theorem targets, and note-layer caveat boundaries.
\end{itemize}

\subsection*{Flagship stubs}

The flagship stub layer is architectural rather than evidentiary. \texttt{ReferenceFlagshipStub} and \texttt{CompletionPartnerFlagshipStub} bundle the abstract ingredients needed for later flagship-level theorem work, but they do not provide concrete empirical witness theorems for the paper's reference-positive or completion-partner cases.

\subsection*{Proof-gap boundary}

The proof-gap layer remains explicit. It records that the current Lean track formalizes abstract definitions, criterion statements, nearby-support distinction vocabulary, and stub interfaces, while leaving open stronger theorems about empirical verdicts, nearby-support synthesis, bundle/dual upgrade claims, indexed versus nonindexed classification, and local-chart consequences. The nearby-positive caveat and the local-chart nonreplacement boundary remain note-layer items there rather than theorem claims.
 \section{Exact Tables}
\label{app:exact-tables}

This appendix collects compact exact tables that support the main text without interrupting the main narrative flow.

\begin{table}[htbp]
\centering
\small
\begin{tabularx}{\linewidth}{@{}l c c c X@{}}
\toprule
Candidate class & Core readout & Core closure & Plateau size & Protocol / verdict role \\
\midrule
Intrinsic candidate & 6/6 & 6/6 & 4 & honest nonindexed baseline; main verdict route \\
Sectioned candidate & 6/6 & 0/6 & 0 & comparison-only; no displacement \\
Bundle-section-fiber candidate & 6/6 & 0/6 & 0 & comparison-only; no displacement \\
Protocol-indexed candidate & 0/6 & 0/6 & 0 & not a decisive supported route \\
\bottomrule
\end{tabularx}
\caption{Core candidate comparison matrix on the supported fixed-support surface.}
\label{tab:appendix-core-candidate-matrix}
\end{table}

\begin{table}[htbp]
\centering
\small
\begin{tabularx}{\linewidth}{@{}>{\raggedright\arraybackslash}p{0.22\linewidth} c c c c >{\raggedright\arraybackslash}p{0.22\linewidth}@{}}
\toprule
Candidate class & Recoarsenings & Rebuilds & Fragile survivors & Collapses & Interpretation \\
\midrule
Intrinsic candidate & 4 & 4 & 0 & 0 & nontrivial stable plateau \\
Sectioned candidate & 2 & 2 & 2 & 0 & fragile survival without plateau \\
Bundle-section-fiber candidate & 3 & 3 & 2 & 1 & fragile survival without plateau \\
Protocol-indexed candidate & 0 & 0 & 0 & 0 & not a decisive surviving route on this surface \\
\bottomrule
\end{tabularx}
\caption{Second-stage plateau summary by candidate on the supported scope.}
\label{tab:appendix-plateau-summary}
\end{table}

\begin{table}[htbp]
\centering
\small
\begin{tabularx}{\linewidth}{@{}>{\raggedright\arraybackslash}p{0.20\linewidth}>{\raggedright\arraybackslash}p{0.14\linewidth}>{\raggedright\arraybackslash}p{0.20\linewidth}>{\raggedright\arraybackslash}X@{}}
\toprule
Audit layer & Result & Nearby positive count & Interpretive note \\
\midrule
Readout-surface test & anchor-only & 0/4 exact readout support & does not show object absence \\
Carrier-level closure suite & carrier support present & 4/4 & not carrier-empty on control-frozen metrics \\
Detector admissibility audit & detector unlock restricted & 0 admissible unlocks & no new theorem-admissible exact-finite detector \\
Family-level quotient audit & erased-partner support & 4/4 erased-like & family structure present despite readout restriction \\
Subset witness search & not needed & 4/4 skipped & all nearby cases resolved at quotient layer \\
Final synthesis & proceed with caveat & route exhausted & final package verdict unchanged \\
\bottomrule
\end{tabularx}
\caption{Nearby-route closeout summary on the frozen admissible surface.}
\label{tab:appendix-nearby-closeout}
\end{table}
 
\clearpage
\bibliographystyle{plainnat}
\bibliography{references}

\end{document}
